\section{R6 with bounded z/B is also an all-tests admission}
\label{sec:r17-r6-bounded-a}
\begin{theorem}
Retain every original R6 hypothesis and normalize $d=1$: $B,z,y>0$
infinitesimal; $t,h,l,j,r$ positive bounded units; $g$ infinitesimal;
the unchanged physicality, contact, coefficient signs, total scale and
prior exclusions. If $a=z/B$ is bounded, then the entire original
parent is finite tensor SOS at all independent projective tests,
including literal axes/infinities, either sign, and complex states.
This admits the whole stated subdomain. The remaining R6 sources have
$a$ unbounded and $1-r$ infinitesimal; no claim about that endpoint is
made here.
\end{theorem}

Write $e=1-r$. Since $1-g=a(e-H)$, the added condition is equivalent
to $e$ being a positive unit; $a_0=1/e_0$. Denote the real residues
by $e\in(0,1)$, $h\in(0,1)$, $l=1-h$, $t>l$, and $\delta=t-l$.
The exact source formulas, with no rate comparison, give
\[
\begin{gathered}
 n_0={e^2\over2(1-e^2)},\quad
 A_0={t\over2h(1+e)},\quad
 C_0={h+t-el\over2l(1+e)},\\
 D_0={1\over2(1+e)\delta},\quad
 (\alpha/B^2)_0={h+t-el\over2e(1+e)},\quad
 \kappa=(mg)_0={1\over2(1+e)}.
\end{gathered}                       \tag{R6.1}
\]
Every listed quantity is a positive unit; $m$ is unbounded. In
particular $n$ cannot be literally zero on this admitted subdomain.
The unbounded-$a$ endpoint, where $n=0$ may occur, is not erased.

Use the exact full-parent excess identity R5.2--R5.5 with a positive
real cap $K$, retaining $(m-K)R_{mc}^2$. The remainder has bounded
coefficients in the unchanged five WA roots, because $a,t$ and all
remaining event weights are bounded and $mg$ bounded. At zero testing
angles its $\eta,R_1$ block is R5.6 and its radial $T$ pivot is
$(\alpha/B^2)_0$. Its $\xi,R_2$ block is
\[
 \begin{pmatrix}n_0+1/2+K&-K-\kappa\\
                 -K-\kappa&K+2\kappa\end{pmatrix},
\quad\det=K(n_0+1/2)+2\kappa(n_0+1/2)-\kappa^2.     \tag{R6.2}
\]
Since $n_0+1/2=\kappa/(1-e)$, its determinant is positive for every
$K>0$. Thus the complete local remainder has strict residue, and WA10
proves all infinitesimal or zero tests, keeping all hard and normal terms.

For a finite remote pair, use the original four states. After centered
extraction the bounded remainder is, in residue,
\[
 B_{\rm fin}+K M^2+2\kappa vWM,\qquad
 M=vW-X+uvY-uZ,                                      \tag{R6.3}
\]
where $B_{\rm fin}$ is the complete original parent at the source
residues R6.1, with $m$ and $\alpha$ omitted. The restriction to
$M=0$ is obtained by the exact substitution $X=vW+uvY-uZ$.
On $(W,Y,Z)$ its determinant is
\[
 { (1+U)^2(1+V)^2\over\mathcal D}
 \left[(h+t)(etU-V)^2+(h+t-el)
          \{t(h+t-e)U^2V+hUV^2\}\right],             \tag{R6.4}
\]
where $\mathcal D=8h(1-e)(1+e)^3l\delta>0$.
Every coefficient outside the displayed square is strictly positive:
$h+t>1>e$ follows from $t>l$. Hence the determinant is positive
whenever $(U,V)\ne(0,0)$, including both axes. The restriction to
$Z=0$ is positive definite: the original signed core there is the sum
of the three squares
$[J^2+(1+U)VY^2+(1+V)X^2]/2$, and the positive $A,C,n$ packets
force $Y=W=0$ on $M=Z=0$ at every such remote pair.
Thus $B_{\rm fin}|_{M=0}$ is positive definite. A sufficiently large
positive real $K$ makes R6.3 strictly positive definite on the full
four-state space. Use this $K$ in the exact source identity and apply
WA10 around the real testing residues. All coefficients remain bounded;
there is no mixed division or source/test comparison.

For unbounded $u$ and bounded $v$, use the full polynomial chart
$p^2F$, $p=1/u$. Its excess correction contains $p$ and disappears in
residue. The centered $m$ row has residue $vY-Z$. On its kernel the
complete finite remainder on $(X,W,Y)$ is
\[
 n_0(1+V)X^2+C_0(1+V)W^2+A_0(1+V)Y^2
                                      +(D_0+1/2)(W+vX)^2\succ0.\tag{R6.5}
\]
A positive cap gives a strict full residue. For unbounded $v$ and bounded
$u$, use $q^2F$, $q=1/v$. The $m$ row is $W+uY$ and the bounded
correction is $2\kappa W(W+uY)$. On its kernel, the complete residue is
\[
 (n_0+1/2)(1+U)X^2+C_0(1+U)^2Y^2
             +(A_0+1/2)(1+U)Y^2+D_0(uX-Z)^2\succ0.  \tag{R6.6}
\]
Again a sufficiently large real $K$ gives a strict full residue.
These assertions include every finite real value of the other angle,
with either sign or zero.

For both angles unbounded use the exact double homogeneous parent and
core R5.16, now in the five roots $(X,W,Y,P=pZ,Q=qZ)$, without normalizing
$X$. Its complete residue after centered extraction is
\[
 (n_0+1/2)X^2+C_0W^2+A_0Y^2+D_0(X-P)^2
                                  +(K+1/2)(Y-Q)^2\succ0.        \tag{R6.7}
\]
The $X,P$ block has determinant $D_0(n_0+1/2)>0$ and the $Y,Q$ block
has determinant $A_0(K+1/2)>0$. The excess correction has a factor $p$
and disappears in residue. All homogeneous coefficients are bounded.
Literal $p=0,q=0$ are the polynomial chart values, not limiting divisions.

These cases exhaust all tests on the stated source subdomain. The
separated finite lift, exact excess restoration, pure-testing homogeneous
restoration and real/imaginary split are exactly those established in
the R5 join. They retain every original normal and hard block, all signs,
all independent rates and the original source/prior-exclusion intersection.
